Esta dissertação apresenta a formulação e a simulação de uma missão de Rendezvous and Docking or Berthing conduzida por uma espaçonave perseguidora equipada com manipulador robótico, com ênfase no encadeamento coerente das etapas de missão e na consistência entre modelos, referenciais e leis de controle ao longo do perfil de operação. A implementação foi realizada em ambiente MATLAB/Simulink, com modelos e controladores desenvolvidos pelo autor, sem o emprego de simuladores prontos de RVDB ou módulos predefinidos. A fase de longa distância é representada por uma transferência de Hohmann entre órbitas circulares coplanares, utilizada para estabelecer condições iniciais para a aproximação de curta distância; nesta fase, o movimento relativo é descrito pelas equações linearizadas de Hill--Clohessy--Wiltshire e regulado por controle proporcional--derivativo, enquanto a atitude do perseguidor é modelada pela equação de Euler para corpo rígido. Na aproximação final, o modelo dinâmico é comutado para um sistema acoplado espaçonave--manipulador obtido por equações de Lagrange, e um controlador PD de atitude atua para mitigar perturbações induzidas pelo movimento das juntas; o manipulador é comandado por referência de cinemática inversa. O desempenho é avaliado por métricas de erro de posição do efetuador final e por métricas geométricas de alinhamento entre o eixo da ferramenta e a normal da porta de acoplamento, ambas expressas no referencial do manipulador. Os resultados obtidos indicam que a arquitetura proposta permite a transição entre fases com estabilidade e coerência de estados, viabilizando a aproximação controlada e o alinhamento geométrico requerido no cenário estudado, ao mesmo tempo em que evidencia a importância da comutação de modelos e do desacoplamento/compensação de perturbações na fase final. As limitações do estudo incluem o uso de dinâmica relativa linearizada para órbitas circulares coplanares e a ausência de perturbações ambientais, o que restringe a generalização direta para cenários operacionais mais complexos.
